\section{The true point-chart initial restriction in the original analytic germ ring}

Use the actual original analytic germ ring $R$, its entire genuine Taylor
map $T$, and the true point-blowup chart. Write $z_e$ for its original
analytic exceptional coordinate germ, and $E(P)$ for actual evaluation of
the finite complex polynomial $P$ at the original coordinate germs.

\begin{lemma}[Actual analytic polynomial evaluation has its entire formal inclusion]
\label{hir:actual-unshifted-analytic-polynomial-taylor-inclusion}
\uses{hir:actual-formal-coefficients-original-polynomial-approximation}
\lean{Hironaka.HironakaActualAnalyticPointInitialRestriction.actualTaylor_polynomial_coordinate_evaluation}
\leanok
For every actual complex polynomial $P$, $T(E(P))$ is its entire formal
polynomial inclusion, with the original coordinate dimension retained.
\end{lemma}
\begin{proof}
The actual Taylor coordinate dictionary sends every original coordinate
germ to the same formal variable. Constants are preserved by the original
complex algebra homomorphism. Polynomial induction through addition and
multiplication then proves the entire identity. No convergence or chosen
polynomial representative certificate is an input.
\end{proof}

\begin{theorem}[Every genuine analytic normalized quotient retains the entire analytic initial restriction]
\label{hir:actual-analytic-point-normalized-initial-restriction}
\uses{hir:actual-unshifted-analytic-polynomial-taylor-inclusion,hir:actual-point-exceptional-restriction-dehomogenized-polynomial,hir:actual-analytic-blowup-coordinate-pullback-taylor-naturality,hir:actual-formal-extension-whole-original-ideal-contraction}
\lean{Hironaka.HironakaActualAnalyticPointInitialRestriction.actual_analytic_point_quotient_initial_restriction}
\leanok
Suppose the actual original germ $g$ lies in the whole point ideal power
$\mathfrak m^d$ and the original analytic quotient $H$ satisfies the true
product identity $z_e^dH=\pi^*g$. Then
\[
 H-E\bigl(\delta_{e,b}P_d(Tg)\bigr)\in(z_e)
\]
in the entire original analytic germ ring.
\end{theorem}
\begin{proof}
Actual whole Taylor naturality carries the true product identity to
$x_e^dT(H)=\pi^*T(g)$. The proved whole point-power image puts $T(g)$
in the actual required formal point power. Genuine exceptional quotient
uniqueness identifies $T(H)$ with that same actual normalized formal
member. The proved whole dehomogenized restriction therefore puts
$T(H)-\delta_{e,b}P_d(Tg)$ in the entire formal ideal $(x_e)$.
The preceding polynomial identity makes this the Taylor image of the
displayed original analytic difference. The genuine coordinate dictionary
makes $(x_e)$ exactly the extension of the whole analytic principal ideal
$(z_e)$. Entire original ideal contraction consequently puts the original
difference in $(z_e)$. The contraction uses the whole original ideal, not
only its finite jets, and never assumes all formal coefficients converge.
\end{proof}

\begin{theorem}[Only the original member and true point-power membership produce the analytic quotient and initial restriction]
\label{hir:actual-analytic-point-initial-restriction-producer}
\uses{hir:actual-analytic-point-normalized-initial-restriction,hir:actual-analytic-blowup-exceptional-quotient-formal-dictionary}
\lean{Hironaka.HironakaActualAnalyticPointInitialRestriction.exists_actual_analytic_point_initial_quotient}
\leanok
Given only the actual original $g\in\mathfrak m^d$ and true point chart,
internally construct an original analytic quotient $H$ satisfying both
$z_e^dH=\pi^*g$ and the entire analytic initial-restriction conclusion above.
\end{theorem}
\begin{proof}
The genuine original analytic exceptional-division theorem internally
produces $H$ and its literal product identity from the stated whole
point-power membership. Apply the preceding original analytic restriction
theorem to this internally constructed same quotient. No quotient,
restriction, ideal equality or output dictionary is supplied by the caller.
\end{proof}

This supplies actual analytic input for the point-centre Hilbert--Samuel
comparison. The all-order inequality and its general permissible-centre
extension remain open.
